\documentclass[11pt]{article}
\usepackage[a4paper,margin=27mm]{geometry}
\usepackage{amsmath,amssymb,mathtools}
\usepackage{fontspec,unicode-math}
\setmainfont{Latin Modern Roman}
\setmathfont{Latin Modern Math}
\usepackage{xurl,hyperref}
\hypersetup{colorlinks=true,urlcolor=blue,linkcolor=blue}
\newcommand{\sref}[2]{\hyperlink{source:#1:#2}{[S#2]}}
\newcommand{\cataloguescope}[1]{\paragraph{Mathematical question.} #1}
\begin{document}
\section{Does every simple first-order theory have stable forking}\label{statement}

\subsection{Definitions and mathematical statement}

\paragraph{Shared notation.}$\mathbb N=\{1,2,\ldots\}$, and $\mathbb Z,\mathbb Q,\mathbb R,\mathbb C$ denote the integers, rationals, reals and complex numbers. Any other conventions are specified locally.


Work in a sufficiently saturated model of a complete first-order theory $T$. A formula $\varphi(x,y)$ is stable if no tuples $(a_i,b_i)_{i\ge0}$ satisfy $\varphi(a_i,b_j)$ exactly when $i<j$. The theory is simple if no formula has the tree property: there are no $k\ge2$ and parameters $b_\eta$ ($\eta\in\mathbb N^{<\mathbb N}$) for which each branch gives a consistent set $\{\varphi(x,b_{\sigma\upharpoonright n}):n\ge0\}$ while each family of immediate successors is $k$-inconsistent. A formula $\psi(x,d)$ divides over $C$ if some $C$-indiscernible sequence $(d_i)$ with $d_0=d$ makes $\{\psi(x,d_i)\}$ inconsistent; it forks over $C$ if it implies a finite disjunction of dividing formulas. A type forks if it contains a forking formula.
\paragraph{Statement recorded in the supplied source.}
In every simple theory, if $\operatorname{tp}(a/Cb)$ forks over $C$, must it contain a formula $\varphi(x,\bar b)$ which forks over $C$ and whose parameter-free formula $\varphi(x,y)$ is stable? This is the historical stable forking conjecture. The cited 2026 paper states a counterexample in Theorem 1.3. \sref{statement}{2}
\subsection{Short English statement}
Must forking in a simple theory always be witnessed by a stable formula? The cited paper claims a negative answer to this historical question.
\subsection{Sources}
\begin{description}
\item[\hypertarget{source:statement:1}{S1}] ULAM.AI and the UnsolvedMath Contributors, UnsolvedMath: A Curated Collection of Open Mathematics Problems, catalogue entry. CC BY 4.0. \url{https://huggingface.co/datasets/ulamai/UnsolvedMath}.
\item[\hypertarget{source:statement:2}{S2}] James Freitag and Scott Mutchnik, A Counterexample to the Stable Forking Conjecture, arXiv:2609.00436v1 (31 August 2026), Definitions 1.1 and Conjecture 1.2, p. 2; Theorem 1.3, p. 3. \url{https://arxiv.org/pdf/2609.00436}.
\end{description}
\label{end:statement}
\end{document}
